\section{Practical limits of imagination}
\label{sec:C06:limits}
World models are valuable when their predictions preserve the distinctions that change a decision, their uncertainty influences control, and real observations correct their state. They are especially attractive when external trials are expensive or risky. Their principal failure is not merely an unattractive generated sample; it is a confidently wrong counterfactual that a planner selects and a memory service later treats as truth. Reliable use therefore requires typed provenance, bounded rollout horizons, real-state verification, and explicit treatment of irreversible effects.

The resource and evidence boundary should remain visible in every claim. A trained actor benefiting from imagined experience is different from online model-predictive control. A visually realistic generator is different from an executable simulator. A latent prediction objective is different from an action-sufficient representation. The strongest comparative studies connect these levels with ablations and verified outcomes instead of allowing one attractive metric to stand in for all of them. Under that standard, world modeling becomes a testable component of agentic intelligence rather than an undefined promise of internal simulation.

\paragraph{A quantitative horizon-selection example.}
The finite-horizon bound can guide an experiment without pretending that its unknown premises have been measured. Suppose, hypothetically, rewards lie in $[0,1]$, uniform reward error is at most $0.01$, and uniform transition total variation is at most $0.001$. At horizon 10, the bound is $0.1+0.045=0.145$; at horizon 40 it is $0.4+0.78=1.18$. These are calculated worst-case return discrepancies under deliberately specified assumptions. They are not empirical errors of Dreamer or a language-model world model. The example shows how horizon changes can dominate a small improvement in average prediction loss. If the uniform transition premise is unavailable, these numbers cannot be used as a deployment certificate.

A more discriminating control experiment estimates errors on planner-induced trajectories and compares short-horizon replanning against long open-loop execution at matched total model calls. Record which candidate states fall outside the training support according to a predeclared diagnostic, then execute a safe subset in the real environment to check that diagnostic. An uncertainty penalty is useful only if it predicts meaningful error on these selected states. Otherwise it can simply encode familiarity with superficial features. The uncertainty model should be calibrated on data held out from both world-model fitting and planner selection, with the remaining distribution-shift limitation made explicit.

Termination prediction deserves separate auditing. If a model incorrectly predicts successful completion, a planner can avoid further checks and inflate imagined return. If it incorrectly predicts failure, it can waste resources repeating an already completed operation. A continuation head is therefore part of the control semantics, not merely an auxiliary reconstruction objective. Test terminal success, terminal failure, recoverable interruption, and unknown remote status as distinct conditions. For long-running agents, the same task can cross multiple sessions, so a process ending is not necessarily an environment terminal state.

Finally, a world model can be evaluated as an evidence-selection mechanism rather than an authoritative simulator. It may predict which real query would resolve uncertainty, then ask the tool for that query. This use can tolerate some prediction error because the final fact comes from the external source. Compare it with random querying, fixed checklists, and a direct policy under the same query budget. The resulting outcome identifies whether imagined interaction improves information gathering, even when the model is insufficiently accurate for open-loop control. Separating these roles prevents an unnecessarily binary judgment that a world model is either universally reliable or useless.
